\section*{Exercises on \S\arabic{exercisegroup}}
\phantomsection
\addcontentsline{toc}{section}{Exercises on \protect\S\arabic{exercisegroup}}

\begin{enumerate}
\def\labelenumi{\arabic{enumi})}
\tightlist
\item
  Let \((f_{\alpha})\) be a set of real functions, defined on an interval \(\mathrm{I} \subset \mathbf{R}\) , each admitting a strict primitive on \(\mathrm{I}\) , and forming a directed set for the relation \(\leqslant\) . Let \(f\) be the upper envelope of the family \((f_{\alpha})\) ; suppose that \(f\) admits a strict primitive on \(\mathrm{I}\) . Show that if \(g_{\alpha}\) (resp. \(g\) ) is the primitive of \(f_{\alpha}\) (resp. \(f\) ) which vanishes at a point \(x_0 \in \mathrm{I}\) , then \(g\) is the upper envelope of the family \((g_{\alpha})\) on the intersection \(\mathrm{I} \cap [x_0, +\infty[\) and its lower envelope on the intersection \(\mathrm{I} \cap ] - \infty\) , \(x_0]\) . (Restricting to the first of these intervals, show that if \(u\) is the upper envelope of \((g_{\alpha})\) on \(\mathrm{I} \cap [x_0, +\infty[\) then one has \(u(x + h) - u(x) \leqslant g(x + h) - g(x)\) for \(h > 0\) ; then prove that, for all \(\alpha\) , one has \(u(x + h) - u(x) \geqslant g_{\alpha}(x + h) - g_{\alpha}(x)\) , and conclude from this last inequality that \(\liminf_{h \to 0} (u(x + h) - u(x))/h \geqslant f(x)\) ; deduce the proposition from this.)
\end{enumerate}

Give an example of an increasing sequence \((f_n)\) of continuous functions on an interval I which are uniformly bounded, but whose upper envelope does not admit a strict primitive on I.

\begin{enumerate}
\def\labelenumi{\arabic{enumi})}
\setcounter{enumi}{1}
\item
  Show that for a function \(\mathbf{f}\) to be a step function on an interval I it is necessary and sufficient that it have only a finite number of points of discontinuity, and be constant on every interval where it is continuous.
\item
  Let \(\mathbf{f}\) be a regulated function on an interval \(I \subset \mathbf{R}\) , taking its values in a complete normed space \(E\) over \(\mathbf{R}\) ; show that for every compact subset \(H\) of \(I\) the set \(\mathbf{f}(H)\) is relatively compact in \(E\) ; give an example where \(\mathbf{f}(H)\) is not closed in \(E\) .
\item
  Give an example of a continuous real function \(f\) on a compact interval \(\mathrm{I} \subset \mathbf{R}\) , such that the composite function \(x \mapsto \operatorname{sgn}(f(x))\) is not regulated on \(\mathrm{I}\) (even though \(\operatorname{sgn}\) is regulated on \(\mathbf{R}\) ).
\item
  Let \(\mathbf{f}\) be a vector function defined on a compact interval \(\mathrm{I} = [a, b] \subset \mathbf{R}\) , taking its values in a complete normed space \(\mathrm{E}\) ; one says that \(\mathbf{f}\) is of bounded variation on \(\mathrm{I}\) if there exists a number \(m > 0\) such that, for every finite strictly increasing sequence \((x_i)_{0 \leqslant i \leqslant n}\) of
\end{enumerate}

points of I such that \(x_0 = a\) and \(x_n = b\) , one has \(\sum_{i=0}^{n-1} \| \mathbf{f}(x_{i+1}) - \mathbf{f}(x_i) \| \leqslant m\) .

\begin{enumerate}
\def\labelenumi{\alph{enumi})}
\item
  Show that \(\mathbf{f}(\mathrm{I})\) is relatively compact in E (argue by contradiction).
\item
  Show that \(\mathbf{f}\) is regulated on I (prove that when \(x\) tends to a point \(x_0 \in \mathrm{I}\) , remaining \(>x_0\) , the function \(\mathbf{f}\) cannot have two different cluster points, and use \(a\) )).
\end{enumerate}

\begin{enumerate}
\def\labelenumi{\arabic{enumi})}
\setcounter{enumi}{5}
\tightlist
\item
  For a function \(\mathbf{f}\) , with values in a complete normed space \(E\) , and defined on an open interval \(I \subset R\) , to be equal to a regulated function on \(I\) at all points of the complement of a countable subset of \(I\) , it is necessary and sufficient that it satisfy the following condition: for every \(x \in I\) and every \(\varepsilon > 0\) there exist a number \(h > 0\) and two elements \(a\) , \(b\) of \(E\) such that one has \(\| \mathbf{f}(y) - a \| \leqslant \varepsilon\) for every \(y \in [x, x + h]\) , except for at most a countably infinite number of points of this interval, and \(\| \mathbf{f}(z) - b \| \leqslant \varepsilon\) for all \(z \in [x - h, x]\) except for at most a countably infinite number of points of this interval.
\end{enumerate}

* 7) Show that the function equal to \(\sin(1/x)\) for \(x \neq 0\) and to 0 for \(x = 0\) , admits a strict primitive on \(\mathbf{R}\) (remark that \(x^2 \sin(1/x)\) admits a derivative at every point).

Deduce from this that if \(g(x,u,v)\) is a polynomial in u, v, its coefficients being continuous functions of x on an interval I containing 0, then the function equal to \(g(x,\sin1/x,\cos1/x)\) for \(x\neq0\) , and to a suitable value \(\alpha\) (to be determined) for x=0, admits a strict primitive on I; give an example where \(\alpha\neq g(0,0,0)\) .

* 8) Show that there exists a continuous function on \([-1, +1]\) admitting a finite derivative at every point of this interval, this derivative being equal to \(\sin\left(\frac{1}{\sin 1/x}\right)\) at points \(x\) different from \(1/n\pi\) ( \(n\) an integer \(\neq 0\) ) and from 0. (On a neighbourhood of \(x = 1/n\pi\) make the change of variable \(x = \frac{1}{n\pi + \operatorname{Arc} \sin u}\) and use exerc. 7; by means of the same change of variable, show that there exists a constant \(a > 0\) independent of \(n\) such that

\[
\left| \int_ {2 / (2 n + 1) \pi} ^ {2 / (2 n - 1) \pi} \sin \left(\frac {1}{\sin \frac {1}{x}}\right) d x \right| \leqslant \frac {a}{n ^ {3}};
\]

and deduce that if one puts

\[
g (x) = \lim _ {\varepsilon \rightarrow 0} \int_ {\varepsilon} ^ {x} \sin \left(\frac {1}{\sin \frac {1}{t}}\right) d t,
\]

then g has a derivative equal to 0 at the point x = 0.) *

\begin{enumerate}
\def\labelenumi{\arabic{enumi})}
\setcounter{enumi}{8}
\item
  Let \(\mathbf{f}\) be a regulated function on a compact interval \(\mathrm{I} = [a, b]\) . Show that for every \(\varepsilon > 0\) there exists a continuous function \(\mathbf{g}\) on \(\mathrm{I}\) such that \(\int_{a}^{b} \| \mathbf{f}(t) - \mathbf{g}(t) \| dt \leqslant \varepsilon\) (reduce to the case where \(\mathbf{f}\) is a step function). Deduce that there exists a polynomial \(\mathbf{h}\) (with coefficients in E) such that \(\int_{a}^{b} \| \mathbf{f}(t) - \mathbf{h}(t) \| dt \leqslant \varepsilon\) .
\item
  Let \(\mathbf{f}\) be a regulated function on \([a, b]\) , taking values in E, let \(\mathbf{g}\) be a regulated function on \([a, c]\) ( \(c > b\) ), taking its values in F, and let \((\mathbf{x}, \mathbf{y}) \mapsto [\mathbf{x}, \mathbf{y}]\) be a continuous bilinear map from \(\mathrm{E} \times \mathrm{F}\) into G (E, F, G being complete normed spaces). Show that
\end{enumerate}

\[
\lim _ {h \rightarrow 0, h > 0} \int_ {a} ^ {b} [ \mathbf {f} (t). \mathbf {g} (t + h) ] d t = \int_ {a} ^ {b} [ \mathbf {f} (t). \mathbf {g} (t) ] d t
\]

(reduce to the case where \(\mathbf{f}\) is a step function).

\begin{enumerate}
\def\labelenumi{\arabic{enumi})}
\setcounter{enumi}{10}
\tightlist
\item
  With the same hypotheses as in exerc. 10 show that for every \(\varepsilon > 0\) there exists a number \(\rho > 0\) such that for every subdivision of \([a, b]\) into intervals \([x_i, x_{i+1}]\) of length \(\leqslant \rho\) ( \(0 \leqslant i \leqslant n-1\) ) one has
\end{enumerate}

\[
\left\| \int_ {a} ^ {b} [ \mathbf {f} (t). \mathbf {g} (t) ] d t - \sum_ {i = 0} ^ {n - 1} [ \mathbf {f} (u _ {i}). \mathbf {g} (v _ {i}) ] (x _ {i + 1} - x _ {i}) \right\| \leqslant \varepsilon
\]

for any choice, for each index i, of points \(u_{i}\) , \(v_{i}\) in \([x_{i}, x_{i+1}]\) (reduce to the case where f and g are step functions).

¶ 12) One says that a sequence \((x_{n})\) of real numbers in the interval \([0, 1]\) is uniformly distributed in this interval if

\[
\lim _ {n \to \infty} \frac {\nu_ {n} (\alpha , \beta)}{n} = \beta - \alpha\tag{1}
\]

for every pair of numbers \(\alpha\) , \(\beta\) such that \(0 \leqslant \alpha \leqslant \beta \leqslant 1\) , where \(\nu_{n}(\alpha, \beta)\) denotes the number of indices \(i\) such that \(1 \leqslant i \leqslant n\) and \(\alpha \leqslant x_{i} \leqslant \beta\) .

Show that if the sequence \((x_{n})\) is uniformly distributed and if \(\mathbf{f}\) is a regulated function on \([0, 1]\) , one has

\[
\lim _ {n \rightarrow \infty} \frac {1}{n} \sum_ {i = 1} ^ {n} \mathbf {f} (x _ {i}) = \int_ {0} ^ {1} \mathbf {f} (t) d t\tag{2}
\]

(reduce to the case where \(\mathbf{f}\) is a step function). Converse.

Show that for the sequence \((x_{n})\) to be uniformly distributed it is enough that the relation (2) should hold for every real function f belonging to a dense set in the space of real continuous functions on {[}0, 1{]} endowed with the topology of uniform convergence.

¶ 13) Let \(f\) be a real regulated function on a compact interval \([a, b]\) . Put

\[
r (n) = \frac {b - a}{n} \sum_ {k = 1} ^ {n} f \left(a + k \frac {b - a}{n}\right) - \int_ {a} ^ {b} f (t) d t.
\]

\begin{enumerate}
\def\labelenumi{\alph{enumi})}
\tightlist
\item
  Show that if \(f\) is increasing on \([a, b]\) one has
\end{enumerate}

\[
0 \leqslant r (n) \leqslant \frac {b - a}{n} \Bigl (f (b) - f (a) \Bigr).
\]

\begin{enumerate}
\def\labelenumi{\alph{enumi})}
\setcounter{enumi}{1}
\tightlist
\item
  If \(f\) is continuous and admits a regulated bounded right derivative on \([a, b[\) , show that one has \(\lim_{n\to \infty}nr(n) = \frac{b - a}{2} (f(b) - f(a))\) (putting \(x_{k} = a + k\frac{b - a}{n}\) ), remark that one has
\end{enumerate}

\[
r (n) = \sum_ {k = 0} ^ {n - 1} \int_ {x _ {k}} ^ {x _ {k + 1}} \left(f (x _ {k + 1}) - f (t)\right) d t,
\]

and apply prop. 5 of II, p.~57).

\begin{enumerate}
\def\labelenumi{\alph{enumi})}
\setcounter{enumi}{2}
\tightlist
\item
  Give an example of a function \(f\) , increasing and continuous on \([a, b]\) , such that \(nr(n)\) does not tend to \(\frac{b - a}{2} (f(b) - f(a))\) as \(n\) grows indefinitely {[}take for \(f\) the limit of a decreasing sequence \((f_n)\) of increasing functions whose graphs are broken lines such that
\end{enumerate}

\[
f _ {n} \left(a + k \frac {b - a}{2 ^ {n}}\right) = f \left(a + k \frac {b - a}{2 ^ {n}}\right) \quad \text { for } \quad 0 \leqslant k \leqslant 2 ^ {n}
\]

and

\[
(b - a) \sum_ {k = 1} ^ {2 ^ {n}} f _ {n} \left(a + k \frac {b - a}{2 ^ {n}}\right) - 2 ^ {n} \int_ {a} ^ {b} f _ {n} (t) d t \geqslant \frac {3}{4} (b - a) \left(f _ {n} (b) - f _ {n} (a)\right) \Biggr ].
\]

\begin{enumerate}
\def\labelenumi{\arabic{enumi})}
\setcounter{enumi}{13}
\tightlist
\item
  Let \(\mathbf{f}\) be a vector function which is the primitive of a regulated function \(\mathbf{f}'\) on \([a, b]\) , and such that \(\mathbf{f}(a) = \mathbf{f}(b) = 0\) . Show that if \(\mathbf{M}\) is the least upper bound of \(\| \mathbf{f}'(x)\|\) over the set of points of \([a, b]\) where \(\mathbf{f}'\) is continuous, then
\end{enumerate}

\[
\left\| \int_ {a} ^ {b} \mathbf {f} (t) d t \right\| \leqslant \mathrm{M} \frac {(b - a) ^ {2}}{4}.
\]

\begin{enumerate}
\def\labelenumi{\arabic{enumi})}
\setcounter{enumi}{14}
\tightlist
\item
  Let \(\mathbf{f}\) be a continuous real function, strictly increasing on an interval \([0, a]\) , and such that \(f(0) = 0\) ; let \(g\) be its inverse function, defined and strictly increasing on \([0, f(a)]\) ; show that
\end{enumerate}

\[
x y \leqslant \int_ {0} ^ {x} f (t) d t + \int_ {0} ^ {y} g (u) d u
\]

for \(0 \leqslant x \leqslant a\) and \(0 \leqslant y \leqslant f(a)\) , equality occurring only if \(y = f(x)\) (study the variation of \(xy - \int_{0}^{x} f(t) dt\) ) as a function of x (y remaining fixed). Deduce that for \(x \geqslant 0\) , \(y \geqslant 0\) , p \textgreater{} 1, \(p' = p/(p - 1)\) , one has \(xy \leqslant ax^{p} + by^{p'}\) for a \textgreater{} 0, b \textgreater{} 0 and \((pa)^{p'}(p'b)^{p} \geqslant 1\) .

\begin{enumerate}
\def\labelenumi{\arabic{enumi})}
\setcounter{enumi}{15}
\tightlist
\item
  Let \(\mathbf{f}\) be a regulated vector function on \(\mathrm{I} = [a, b] \subset \mathbf{R}\) , let \(\mathbf{u}\) be a primitive of \(\mathbf{f}\) on \(\mathrm{I}\) and \(\mathrm{D}\) a closed convex set containing \(\mathbf{u}(\mathrm{I})\) . Show that if \(g\) is a real monotone function on \(\mathrm{I}\) one has
\end{enumerate}

\[
\int_ {a} ^ {b} \mathbf {f} (t) g (t) d t = (\mathbf {u} (b) - \mathbf {c}) g (b) + (\mathbf {c} - \mathbf {u} (a)) g (a)
\]

where \(\mathbf{c}\) belongs to D (reduce to the case where \(g\) is a monotone step function). Deduce that if \(f\) is a real regulated function on I then there exists a \(c \in \mathrm{I}\) such that

\[
\int_ {a} ^ {b} f (t) g (t) d t = g (a) \int_ {a} ^ {c} f (t) d t + g (b) \int_ {c} ^ {b} f (t) d t
\]

(``the second mean value theorem'').

\begin{enumerate}
\def\labelenumi{\arabic{enumi})}
\setcounter{enumi}{16}
\tightlist
\item
  Let \(g\) be a real function admitting a continuous derivative, and \(\neq 0\) on \([a, x]\) ; if \(f\) is a real function having a regulated \((n + 1)^{th}\) derivative on \([a, x]\) , show that the remainder \(r_n(x)\) in the Taylor expansion of order \(n\) for \(f\) at the point \(a\) can be written
\end{enumerate}

\[
r _ {n} (x) = (g (x) - g (a)) \frac {(x - \xi) ^ {n}}{n !} \frac {f ^ {(n + 1)} (\xi)}{g ^ {\prime} (\xi)}
\]

where \(a < \xi < x\) (use the integral form for \(r_{n}(x)\) ).

\begin{enumerate}
\def\labelenumi{\arabic{enumi})}
\setcounter{enumi}{17}
\tightlist
\item
  Let \(f\) be a finite real function, continuous on an open interval \(I\) . For \(f\) to be convex on \(I\) it is necessary and sufficient that for every \(x \in I\) one has
\end{enumerate}

\[
\operatorname * {l i m s u p} _ {h \to \infty} \left(\frac {1}{2 h} \int_ {x - h} ^ {x + h} f (t) d t - f (x)\right) \geqslant 0
\]

(argue as in exerc. 9 of I, p.~46).

\begin{enumerate}
\def\labelenumi{\arabic{enumi})}
\setcounter{enumi}{18}
\tightlist
\item
  Let \(f\) be a convex function on an interval \(I\) , let \(h\) be a number \(> 0\) , and \(I_h\) the intersection of \(I\) and the intervals \(I + h\) and \(I - h\) ; show that, if \(I_h\) is not empty, the function
\end{enumerate}

\[
g _ {h} (x) = \frac {1}{2 h} \int_ {x - h} ^ {x + h} f (t) d t
\]

is convex on \(\mathrm{I}_h\) ; if \(h < k\) one has \(g_h \leqslant g_k\) . When \(h\) tends to 0 show that \(g_h\) tends uniformly to \(f\) on every compact interval contained in the interior of \(\mathrm{I}\) .

\begin{enumerate}
\def\labelenumi{\arabic{enumi})}
\setcounter{enumi}{19}
\tightlist
\item
  Show that as \(n\) increases indefinitely the polynomial
\end{enumerate}

\[
f _ {n} (x) = \frac {\int_ {0} ^ {x} (1 - t ^ {2}) ^ {n} d t}{\int_ {0} ^ {1} (1 - t ^ {2}) ^ {n} d t}
\]

tends uniformly to -1 on every interval \([-1, -\varepsilon]\) , and tends uniformly to +1 on every interval \([\varepsilon, +1]\) , where \(\varepsilon > 0\) (note that \(\int_{0}^{1}(1-t^{2})^{n}dt \geqslant \int_{0}^{1}(1-t)^{n}dt\) ). Deduce that the polynomial \(g_{n}(x) = \int_{0}^{x} f_{n}(t)dt\) tends uniformly to \(|x|\) on \([-1, +1]\) , which provides a new proof of the Weierstrass theorem (Gen.~Top., X, p.~313, prop. 3).

\begin{enumerate}
\def\labelenumi{\arabic{enumi})}
\setcounter{enumi}{20}
\tightlist
\item
  Let f be a real increasing convex continuous function on an interval \([0, a]\) , and such that \(f(0) = 0\) . If \(a_{1} \geqslant a_{2} \geqslant \cdots \geqslant a_{n} \geqslant 0\) is a finite decreasing sequence of points in \([0, a]\) , show that one has
\end{enumerate}

\[
f \left(a _ {1}\right) - f \left(a _ {2}\right) + \dots + (- 1) ^ {n - 1} f \left(a _ {n}\right) \geqslant f \left(a _ {1} - a _ {2} + \dots + (- 1) ^ {n - 1} a _ {n}\right).
\]

(One can restrict oneself to the case where \(n = 2m\) is even; remark that for \(1 \leqslant j \leqslant m\) one has

\[
\int_ {\alpha} ^ {\beta} f ^ {\prime} (t) d t \leqslant \int_ {a _ {2 j}} ^ {a _ {2 j - 1}} f ^ {\prime} (t) d t
\]

where \(\alpha = a_{2j+1} - a_{2j+2} + \cdots + a_{2m-1} - a_{2m}\) and \(\beta = \alpha + (a_{2j-1} - a_{2j})\) .

\begin{enumerate}
\def\labelenumi{\arabic{enumi})}
\setcounter{enumi}{21}
\tightlist
\item
  Let f be a continuous increasing real function and \(\geqslant 0\) on the interval {[}0, 1{]}.
\end{enumerate}

\begin{enumerate}
\def\labelenumi{\alph{enumi})}
\item
  Show that there exists a convex function \(g \geqslant 0\) on {[}0,1{]} such that \(g \leqslant f\) and \(\int_0^1 g(t)dt \geqslant \frac{1}{2}\int_0^1 f(t)dt\) (cf.~I, p.~48, exerc. 23).
\item
  Show that there exists a convex function \(h\) on {[}0, 1{]} such that \(h \geqslant f\) and that
\end{enumerate}

\[
\int_ {0} ^ {1} h (t) d t \leqslant 2 \int_ {0} ^ {1} f (t) d t.
\]

(For every \(a\) such that \(0 < a \leqslant 1\) , let \(f_a\) be the function equal to \(f\) for \(0 \leqslant t \leqslant a\) and to \(f(a)\) for \(a \leqslant t \leqslant 1\) ; let A be the set of \(a \in [0,1]\) for which there exists a convex function \(h_a\) on \([0,1]\) such that \(h_a \geqslant f_a\) and \(\int_0^1 h_a(t)dt \leqslant 2\int_0^1 f_a(t)dt\) . Show that the least upper bound \(b\) of A again belongs to A, using exerc. 1 of I, p.~45. Then prove that \(f_{b} = f\) , arguing by contradiction. For this, reduce to proving the following result: if \(\varphi\) is continuous and increasing on \([0, 1]\) , not constant, and such that \(\varphi(0) = 0\) , then there is a point \(c \in ]0, 1[\) such that \(\varphi(c) > 0\) and such that the linear function \(\psi\) such that \(\psi(c) = \varphi(c)\) and

\[
\psi (2 c - 1) = 0
\]

satisfies the relation \(\psi(t) \geqslant \varphi(t)\) for \(\max(0, 2c - 1) \leqslant t \leqslant c\) .

\begin{enumerate}
\def\labelenumi{\arabic{enumi})}
\setcounter{enumi}{22}
\tightlist
\item
  Let \(\mathbf{f}\) be a vector function, continuously differentiable on an interval \([a, b] \subset \mathbf{R}\) , with values in a complete normed space.
\end{enumerate}

\begin{enumerate}
\def\labelenumi{\alph{enumi})}
\tightlist
\item
  Show that for \(a \leqslant t \leqslant b\) one has
\end{enumerate}

\[
\mathbf {f} (t) = \frac {1}{b - a} \int_ {a} ^ {b} \mathbf {f} (x) d x + \int_ {a} ^ {b} \frac {x - a}{b - a} \mathbf {f} ^ {\prime} (x) d x + \int_ {t} ^ {b} \frac {x - b}{b - a} \mathbf {f} ^ {\prime} (x) d x.
\]

\begin{enumerate}
\def\labelenumi{\alph{enumi})}
\setcounter{enumi}{1}
\tightlist
\item
  Deduce the inequalities
\end{enumerate}

\[
\| \mathbf {f} (t) \| \leqslant \frac {1}{b - a} \int_ {a} ^ {b} \| \mathbf {f} (x) \| d x + \int_ {a} ^ {b} \left\| \mathbf {f} ^ {\prime} (x) \right\| d x
\]

and

\[
\left\| \mathbf {f} \left(\frac {1}{2} (a + b)\right) \right\| \leqslant \frac {1}{b - a} \int_ {a} ^ {b} \| \mathbf {f} (x) \| d x + \frac {1}{2} \int_ {a} ^ {b} \left\| \mathbf {f} ^ {\prime} (x) \right\| d x.
\]

\setcounter{exercisegroup}{1}
\refstepcounter{exercisegroup}
